\documentclass[10pt,a4paper]{amsart}
\usepackage[margin=19mm]{geometry}
\usepackage{amsmath,amssymb,mathtools}
\usepackage[scr=rsfs,cal=euler]{mathalfa}
\usepackage{xfrac}
\usepackage[unicode,hidelinks,bookmarks=false]{hyperref}
\newcommand{\ie}{i.e.\ }
\newcommand{\cf}{cf.\ }
\newcommand{\resp}{respectively\ }
\newcommand{\Subsection}{\subsection}
\providecommand{\nobreakdash}{\nobreak\hbox{-}}
\setlength{\emergencystretch}{2em}
\multlinegap0pt
\usepackage{amssymb}
\usepackage{xcolor}
\usepackage{algorithm}\usepackage{algpseudocode}
\newtheorem{proposition}{Proposition}
\title{Evaluation codes from linear systems of conics}
\author{Barbara Gatti, G\'abor Korchm\'aros, Gioia Schulte}
\date{Source: arXiv:2605.11187v1 (CC BY 4.0)}
\begin{document}
\maketitle

\noindent\emph{Source and licence.} Evaluation codes from linear systems of conics. Version arXiv:2605.11187v1 (CC BY 4.0), distributed by arXiv under the Creative Commons Attribution 4.0 International licence, http://creativecommons.org/licenses/by/4.0/, as recorded in the arXiv licence metadata for this exact version and shown on the abstract page https://arxiv.org/abs/2605.11187v1. Full licence notice: \texttt{licenses/arxiv-2605.11187-CC-BY-4.0.txt}.\par\smallskip

\noindent\textit{Editorial scope.} Counted unit: the article's proposition pro12082025 (source0.tex lines 494--496). The article's complete original proof of it is delivered with it (source0.tex lines 497--610, one \texttt{proof} environment of 114 lines). No mathematics is written, repaired or re-derived: the statement and the proof are the article's own lines, in the article's own notation, and no retained line is substituted or re-flowed.  Retained uncounted as the source-local context the proof needs: lines 469--474; lines 476--481.  Nothing was omitted from any retained span, so the package carries no omission record.  No retained line contains a citation, so the wrapper carries no bibliography and no bibliography entry.  No retained line contains a figure environment, an image-inclusion command or drawing code, so no image is delivered and the package declares no asset dependency.  Editorial formatting only, in addition: an AMS \texttt{amsart} preamble that restates the article's own declarations and macros used by the retained lines, namely the environments \texttt{proposition} on separate counters, together with the standard packages those lines need.  The retained environments print in this document as Proposition 1 in order of appearance, whereas the article prints the same items with the numbering of its own full document.  The complete native source of the article as distributed in the arXiv e-print for this exact version is bundled unchanged as \texttt{sources/200344/source0.tex}.
\begin{equation}
\label{eq09082025}
\begin{array}{llll}H(X,V)=&&(Va_{22}+a_{12})X^2+(V^2a_{12}+Va_{23}+\bar{v}^2a_{12}+\bar{v}a_{23}+a_{13})X+\\
&&(\bar{v}a_{23}+a_{13})V^2+Va_{33}+\bar{v}^3a_{23}+\bar{v}^2a_{13}.
\end{array}
\end{equation}

\begin{equation}
\label{eq09082025A}
\begin{array}{llll}H(X,V)=&&(a_{12}X+\bar{v}a_{23}+a_{13})V^2+(X^2a_{22}+Xa_{23}+a_{33})V+\\
&& (X+\bar{v}^2)(a_{12}X+\bar{v}a_{23}+a_{13}).
\end{array}
\end{equation}

\begin{proposition}
\label{pro12082025} Assume that $a_{12}\ne 0$ and $a_{22}\ne 0$. Then $C$ is degenerate if and only if $\mathcal{H}$ is reducible. Moreover, if $\mathcal{H}$ is reducible then it splits into three linear components.
\end{proposition}

\begin{proof}
Let $\ell$ be a line and suppose $\ell$ to be a linear component of $\mathcal{H}$. Then $\ell$ passes through one of the points $V_\infty$, $X_\infty$ and $Q_\infty$.

(i)
If $V_\infty\in \ell$, let $\ell_V=\ell$. Then $\ell_V$ has equation $X-x=0$ for some $x\in \mathbb{F}$. Therefore, $\ell_V$ is a component of $\mathcal{H}$ if and only if $x$ is a solution of the system in the indeterminate $Z$ arising from (\ref{eq09082025A})
\begin{equation}
\label{eq09082025AA}
\begin{cases}
a_{12}Z+\bar{v}a_{23}+a_{13}=0;\\
a_{22}Z^2+a_{23}Z+a_{33}=0;\\
(Z+\bar{v}^2)(a_{12}Z+\bar{v}a_{23}+a_{13})=0.
\end{cases}
\end{equation}
 The first equation in (\ref{eq09082025AA}) implies the third one.
 Moreover, eliminating $Z$ from the first two equations gives $U_{12}=0$ where
\begin{equation}
\label{eq12082025}
U_{12}=\bar{v}^2a_{22}a_{23}^2+\bar{v}a_{12}a_{23}^2+a_{12}^2a_{33}+a_{12}a_{13}a_{23}+a_{22}a_{13}^2.
\end{equation}
Since $a_{11}+a_{12}\bar{v}+a_{22}\bar{v}^2=0$, it turns out from (\ref{eq12082025}) that the system of the first two equations has a solution if and only if
$$a_{11}a_{23}^2+a_{12}a_{23}a_{13}+a_{22}a_{13}^2+a_{33}a_{12}^2=0,$$
that is, $V_\infty\in \ell_V$ if and only if the conic $C$ is reducible.

(ii) If $Q_\infty \in \ell$, let $\ell_Q=\ell$. Then $\ell_Q$ is parallel to the line  of equation $a_{22}X+a_{12}V=0$. Replacing $a_{22}X+a_{12}V$ by $W$, Equation (\ref{eq09082025A}) of $\mathcal{H}$ becomes
\begin{equation}
\label{eq09082025w}
\begin{array}{llll}L(X,W)=&&(Xa_{12}+\bar{v}a_{23}+a_{13})W^2+a_{12}(X^2a_{22}+Xa_{23}+a_{33})W+\\
&&X^2(a_{12}^3+\bar{v}a_{22}^2a_{23}+a_{11}a_{22}a_{33}+a_{22}^2a_{13})+\\
&&a_{12}x(\bar{v}^2a_{12}^2+\bar{v}a_{12}a_{23}+a_{12}a_{13}+a_{22}a_{33})+\\
&& \bar{v}^2a_{12}^2(\bar{v}a_{23}+a_{13}).
\end{array}
\end{equation}
Also, the equation of $\ell_Q$ becomes $X=w$ for some $w\in \mathbb{F}$.
Therefore, $\ell_Q$ is component of $\mathcal{H}$ if and only if $w$ is a solution of the system in the indeterminate $Z$ arising from (\ref{eq09082025w})
\begin{equation}
\label{eq09082025WW}
\begin{cases}
a_{12}Z+\bar{v}a_{23}+a_{13}=0;\\
a_{12}(a_{22}Z^2+a_{23}Z+a_{33})=0;\\
(a_{12}^3+\bar{v}a_{22}^2a_{23}+a_{11}a_{22}a_{33}+a_{22}^2a_{13})
Z^2+\\
a_{12}(\bar{v}^2a_{12}^2+\bar{v}a_{12}a_{23}+a_{12}a_{13})Z+a_{11}a_{22}a_{33}+
a_{12}^2(\bar{v}^3a_{23}+\bar{v}a_{13})=0.
\end{cases}
\end{equation}
The system of the first two equations has a solution if and only if
$Z_{12}=0$ where
\begin{equation}
\label{eq12082025B}
Z_{12}=\bar{v}^2a_{22}a_{23}^2+\bar{v}a_{12}a_{23}^2+a_{12}^2a_{33}+a_{12}a_{13}a_{23}+a_{22}a_{13}^2,
\end{equation}
and that of the first and the third ones if and only if $Z_{13}=0$ where
\begin{equation}
\label{eq12082025BB}
Z_{13}=(\bar{v}a_{23}+a_{13})Z_{12}.
\end{equation}
Therefore, $\ell$ is a component of $\mathcal{H}$ through $Q_\infty$ if and only if $Z_{12}=0$. Since $Z_{12}=a_{22}U_{12}$, it turns out that
$\ell_Q$ is a component of $\mathcal{H}$ if and only if so is $\ell_V$. Moreover, this is the case if and only if $C$ is reducible.

Before dealing with the last case $X_\infty\in \ell$, we rewrite the above arguments starting off with Equation (\ref{eq09082025}).
Replacing $a_{22}X+a_{12}V$ by $W$, Equation (\ref{eq09082025}) of $\mathcal{H}$ becomes
\begin{equation*}
\label{eq09082025w1}
\begin{array}{llll}L(V,W)=&&(Va_{22}+a_{12})W^2+a_{22}(V^2a_{12}+Va_{23}+\bar{v}^2a_{12}+\bar{v}a_{23}+a_{13})W+\\
&&V^2(\bar{v}a_{22}^2a_{23} + a_{12}^3 + a_{12}a_{22}a_{23} + a_{22}^2a_{13}) +  \\ && V(\bar{v}^2a_{12}^2a_{22} +\bar{v}a_{12}a_{22}a_{23} +a_{12}a_{22}a_{13} + a_{22}^2a_{33}) + \bar{v}^3a_{22}^2a_{23} + \bar{v}^2a_{22}^2a_{13}.
\end{array}
\end{equation*}
Also, the equation of $\ell_X$ becomes $V=v$ for some $v\in \mathbb{F}$. 
The formulas analogous to (\ref{eq09082025WW}), (\ref{eq12082025B}), and (\ref{eq12082025BB}) are
\begin{equation*}
\label{eq09082025WR}
\begin{cases}
Za_{22}+a_{12}=0;\\
a_{22}(Z^2a_{12}+Za_{23}+\bar{v}^2a_{12}+\bar{v}a_{23}+a_{13})=0;\\
Z^2\bar{v}a_{22}^2a_{23} + Z^2a_{12}^3 + Z^2a_{12}a_{22}a_{23} + Z^2a_{22}^2a_{13} + Z\bar{v}^2a_{12}^2a_{22} + \\
Z\bar{v}a_{12}a_{22}a_{23} +Za_{12}a_{22}a_{13} + Za_{22}^2a_{33} + \bar{v}^3a_{22}^2a_{23} + \bar{v}^2a_{22}^2a_{13}=0;
\end{cases}
\end{equation*}
$Q_{12}=0$ where
\begin{equation*}
\label{eq12082025BQ}
Q_{12}=\bar{v}^2a_{12}a_{22}^3 + \bar{v}a_{22}^3a_{23}+ a_{12}^3a_{22} + a_{12}a_{22}^2a_{23} + a_{22}^3a_{13};
\end{equation*}
$Q_{13}=0$ where
\begin{equation*}
\label{eq12082025BBQ}
Q_{13}=\bar{v}^3a_{22}^4a_{23}+ \bar{v}^2a_{12}^3a_{22}^2 + \bar{v}^2a_{22}^4a_{13}+ a_{12}^5 + a_{12}^3a_{22}a_{23} + a_{12}a_{22}^3a_{33}.
\end{equation*}
Therefore, $\ell_Q$ is a component of $\mathcal{H}$ if and only if $Q_{12}=0$ and $Q_{13}=0$. If this is the case then $C$ is reducible.

(iii) If $X_\infty\in \ell$, let $\ell_X=\ell$. Then $\ell_X$ has equation $V=v$ for some $v\in \mathbb{F}$. Therefore, $\ell_X$ is a component of $\mathcal{H}$ if and only if $v$ is a solution of the system in the indeterminate $Z$ arising from (\ref{eq09082025})
\begin{equation*}
\label{eq09082025C}
\begin{cases}Za_{22}+a_{12}=0; \\
Z^2a_{12}+Za_{23}+\bar{v}^2a_{12}+\bar{v}a_{23}+a_{13}=0;\\
Z^2(\bar{v}a_{23}+a_{13})+Za_{33}+\bar{v}^3a_{23}+\bar{v}^2a_{13}=0.
\end{cases}
\end{equation*}
The system of the first two equations has a solution if and only if
$R_{12}=0$ where
\begin{equation*}
\label{eq12082025D}
R_{12}=\bar{v}^2 a_{12}a_{22}^2+\bar{v}a_{22}^2a_{23}+a_{12}^3+a_{12}a_{22}a_{23}+a_{22}^2a_{13},
\end{equation*}
and that of the first and the third ones if and only if $R_{13}=0$ where
\begin{equation*}
\label{eq12082025DD}
R_{13}=\bar{v}^3 a_{22}^2a_{23}+\bar{v}^2a_{22}^2a_{13}+\bar{v}a_{12}^2a_{23}+a_{12}a_{22}a_{33}+a_{12}^2a_{13}.
\end{equation*}
By a straightforward computation,
$${\mbox{$Q_{12}=a_{22}R_{12}$ and $Q_{13}=a_{22}^2R_{13}+a_{12}^2R_{12}$}}.$$

It turns out that $\ell_X$ is a component of $\mathcal{H}$ if and only if so is $\ell_Q$ (and hence $\ell_V$). Therefore, $C$ is reducible if and only if $\mathcal{H}$ is reducible, and if $\mathcal{H}$ is reducible then it splits into three lines.
\end{proof}

\end{document}
